Die Kolben gleiten reibungsfrei, ihre Gewichtskräfte und der Schweredruck in der Flüssigkeit werden vernachlässigt; angegeben wird der Überdruck gegenüber der Umgebungsluft.
\begin{abcliste}
\abc Der Körper drückt mit seiner Gewichtskraft auf den Kolben $A_1$. Im Gleichgewicht herrscht in der Flüssigkeit zwischen den Kolben der Druck
\begin{align}
 p &= \pkoF \\
   &= \frac{\mko\cdot\gE}{\Aea} \\
   &= \pko \approx \result{\pkoP{3}{2}}
\end{align}
also rund \pkobP{0}{2} Überdruck.
\abc Die Flüssigkeit drückt den Kolben $A_2$ mit $p\,A_2$ nach links und den Kolben $A_3$ mit $p\,A_3$ nach rechts. Da die beiden Kolben durch den Draht verbunden sind, bilden sie zusammen einen Körper, auf den die Flüssigkeit insgesamt die Kraft $p\left(A_2-A_3\right)$ nach links ausübt. Die Kraft $F$ nach rechts muss diese ausgleichen:
\begin{align}
 F &= p\left(A_2-A_3\right) = \FkoF \\
   &= \frac{\mko\cdot\gE}{\Aea}\cdot\left(\Aeb-\Aec\right) \\
   &= \Fko \approx \result{\FkoP{0}{2}}
\end{align}
\end{abcliste}
